\documentclass[10pt,a4paper]{amsart}
\usepackage[margin=19mm]{geometry}
\usepackage{amsmath,amssymb,mathtools}
\usepackage[scr=rsfs,cal=euler]{mathalfa}
\usepackage{xfrac}
\usepackage[unicode,hidelinks,bookmarks=false]{hyperref}
\newcommand{\ie}{i.e.\ }
\newcommand{\cf}{cf.\ }
\newcommand{\resp}{respectively\ }
\newcommand{\Subsection}{\subsection}
\providecommand{\nobreakdash}{\nobreak\hbox{-}}
\setlength{\emergencystretch}{2em}
\multlinegap0pt
\usepackage{amssymb}
\usepackage[shortlabels]{enumitem}
\usepackage{xcolor}
\usepackage{float}
\newtheorem{proposition}{Proposition}
\newtheorem{lemma}{Lemma}
\newtheorem{theorem}{Theorem}
\newcommand{\F}{\mathbb F}
\newcommand{\Z}{\mathbb Z}
\title{The combinatorial structure and value distributions of plateaued functions}
\author{Lukas K\"olsch, Alexandr Polujan}
\date{Source: arXiv:2410.00611v3 (CC BY 4.0)}
\begin{document}
\maketitle

\noindent\emph{Source and licence.} The combinatorial structure and value distributions of plateaued functions. Version arXiv:2410.00611v3 (CC BY 4.0), distributed by arXiv under the Creative Commons Attribution 4.0 International licence, http://creativecommons.org/licenses/by/4.0/, as recorded in the arXiv licence metadata for this exact version and shown on the abstract page https://arxiv.org/abs/2410.00611v3. Full licence notice: \texttt{licenses/arxiv-2410.00611-CC-BY-4.0.txt}.\par\smallskip

\noindent\textit{Editorial scope.} Counted unit: the article's theorem thm:platdto1 (source0.tex lines 374--376). The article's complete original proof of it is delivered with it (source0.tex lines 377--412, one \texttt{proof} environment of 36 lines). No mathematics is written, repaired or re-derived: the statement and the proof are the article's own lines, in the article's own notation, and no retained line is substituted or re-flowed.  Retained uncounted as the source-local context the proof needs: lines 180--182; lines 204--207; lines 370--372.  Nothing was omitted from any retained span, so the package carries no omission record.  No retained line contains a citation, so the wrapper carries no bibliography and no bibliography entry.  No retained line contains a figure environment, an image-inclusion command or drawing code, so no image is delivered and the package declares no asset dependency.  Editorial formatting only, in addition: an AMS \texttt{amsart} preamble that restates the article's own declarations and macros used by the retained lines, namely the environments \texttt{proposition}, \texttt{lemma}, \texttt{theorem} on separate counters, together with the standard packages those lines need.  The retained environments print in this document as Proposition 1, Lemma 1, Theorem 1 in order of appearance, whereas the article prints the same items with the numbering of its own full document.  The complete native source of the article as distributed in the arXiv e-print for this exact version is bundled unchanged as \texttt{sources/166088/source0.tex}.
\begin{equation} \label{eq:first}
    \sum_{b \in \F_{p}^m}W_F(b,0) = \sum_{x \in \F_{p}^n}\sum_{b \in \F_{p}^m} \zeta_p^{\langle b,F(x)\rangle} = p^m |F^{-1}(0)|.
\end{equation}

\begin{proposition} \label{prop:starting}
Let $F \colon \F_{p}^n \rightarrow \F_{p}^m$ be a function. Then 
\[\sum_{\beta \in \F_{p}^m} |F^{-1}(\beta)|^2 = p^{2n-m}+\frac{1}{p^m} \sum_{b \in \F_p^m\setminus\{0\}} |W_F(b,0)|^2.\]
\end{proposition}

\begin{lemma} \label{lem:podd}
    Let $F \colon \F_p^n \rightarrow \F_p^n$ be a plateaued function that is $d$-to-1 for some $d|p^n-1$, $d>2$, where $F(0)=0$ and $|F^{-1}(0)|=1$. Then $W_F(b,0) \in \Z$ for any $b \in \F_p^n$.
\end{lemma}

\begin{theorem} \label{thm:platdto1}
    Let $F \colon \F_p^n \rightarrow \F_p^n$ be a plateaued function that is $d$-to-1 for some $d|p^n-1$, $d>2$. Then $n$ is even, $d=p^t+1$ for some $t \geq 1$, $t|n/2$ and $F$ has $p^n-1-\frac{p^n-1}{p^t+1}$ bent components and $\frac{p^n-1}{p^t+1}$ components with amplitude $p^{n/2+t}$. The linearity of such a function is thus $\mathcal{L}(F)=p^{n/2+t}$.
\end{theorem}

\begin{proof}
    If $F(x)$ is plateaued then so is $F(x+a)+b$ for any constants $a,b$ and all amplitudes of the component functions remain the same, so we can assume without loss of generality that $F^{-1}(0)=\{0\}$, i.e., $F(0)=0$ and $0$ is the unique element with exactly one preimage.

    Since $F$ is $d$-to-$1$ we have $W_F(b,0) \equiv 1 \pmod{d}$, so $W_F(b,0)-1=d\cdot C$ for some $C \in \Z[\zeta_p]$. In particular, $W_F(b,0) \neq 0$ for all $b \in \F_p^n \setminus \{0\}$. By Lemma~\ref{lem:podd}, we have that $W_F(b,0) \in \Z$, so $W_F(b,0) \in \{-p^{k},p^k\}$.
    
    This implies that $d|p^{k}\pm 1$. 
    Let us first assume that $d\nmid p^t+1$ for all $0<t<n$. Then $W_F(b,0) =  p^{k}$ for some $k$ for any nonzero $b$. But we have (see Eq.~\eqref{eq:first})
    \begin{equation} \label{eq:firstmoment}
        \sum_{b \in \F_p^n \setminus \{0\}}W_F(b,0) =\sum_{b \in \F_p^n}W_F(b,0)-p^n = p^n(|F^{-1}(0)|-1)=0,
    \end{equation}
    yielding a contradiction. We conclude that there is a minimal $0<t<n$ such that $d|p^t+1$.  Assume now that $k\geq n/2$ is another integer such that $W_F(b,0)= \pm p^{k}$ for some $b$. Then $d|p^{k} \pm 1$ and $d|p^t+1$, so in particular $d|p^{\gcd(t,k)}+1$. By the minimality of $t$, we have $\gcd(t,k)=t$. We conclude that the only possible values of $W_F(b,0)$ are $W_F(b,0)= \pm p^{k+it}$ with $i \geq 0$ for some $k \geq n/2$ that is divisible by $t$. Define 
    \[N_i = |\{b\in \F_2^n \colon W_F(b,0)=\pm p^{k+it}\}|\]
    for $i\geq 0$. {Note that for a fixed $i$, only one sign can occur since otherwise $d|p^{k+it}-1$ and $d|p^{k+it}+1$, so $d|2$, contradicting our assumption.} 
    Then we have using Eq.~\eqref{eq:firstmoment}
    \[0=\sum_{b \in \F_p^n \setminus \{0\}}W_F(b,0) = p^{k}(\pm N_0\pm p^tN_1 \pm p^{2t}N_2\pm \dots),\]
    so 
    \begin{equation} \label{eq:1}
        N_0\mp p^tN_1 \mp p^{2t}N_2\mp \dots =0.
    \end{equation}
    Moreover, we have by Proposition~\ref{prop:starting}
    \[    1+d(p^n-1)=\sum_{\beta \in \F_{p}^n} |F^{-1}(\beta)|^2 = p^n + \frac{1}{p^n}\sum_{\beta \in \F_p^n}|W_F(b,0)|^2,
    \]
    so 
    \begin{equation} \label{eq:2}
        (d-1)(p^n-1)=p^{2k-n}(N_0+p^{2t}N_1+p^{4t}N_2+\dots).
    \end{equation}
    Since we have no balanced components, we also have
    \begin{equation} \label{eq:3}
        {N_0}+N_1+N_2+\dots = p^n-1.
    \end{equation}
    We now subtract Eq.~\eqref{eq:3} $(d-1)$-times from Eq.~\eqref{eq:2} and get
    \[(p^{2k-n}+1-d)N_0+(p^{2k-n+2t}+1-d)N_1 + (p^{2k-n+4t}+1-d)N_2+\dots=0.\]
    From this equation, we add Eq.~\eqref{eq:1} $(p^t-1)$-times:
    \[(p^{2k-n}-d+p^t)N_0+(p^{2k-n+2t}+1-d\pm (p^{2t}- p^t))N_1+(p^{2k-n+4t}+1-d\pm(p^{3t}-p^{2t}))N_2+\dots =0.\]
    Note that $d\leq p^t+1$ and $p^{2k-n}\geq 1$, so all coefficients that appear in front of the $N_i$ are non-negative. More specifically, we see that all coefficients are positive, unless $n=2k$, $d=p^t+1$, in which case the coefficients of $N_0$ and $N_1$ are $0$ (when choosing the correct sign in the case of $N_1$), while the others are still positive. We conclude that necessarily $N_i=0$ for all $i \geq 2$ and $n=2k$, $d=p^t+1$. Eqs.~\eqref{eq:2} and~\eqref{eq:3} then yield $N_0+N_1=p^n-1$ and $p^t(p^n-1)=N_0+p^{2t}N_1$, leading to (after solving the two linear equations) $N_0=p^n-1-\frac{p^n-1}{p^t+1}$ and $N_1=\frac{p^n-1}{p^t+1}$. Since $t$ divides $k$ and $2k=n$, we have that $t$ necessarily divides $n/2$.
\end{proof}

\end{document}
